\chapter{שיווק ומכירות}\label{ch:marketing}

גגון מוכרת בשני צעדים: קודם משכנעת ועד בית להשכיר גג, ואחר כך מוכרת את התוצרת מהגג למסעדות ולתושבים. השיווק בנוי לפי הרצף הזה. הוא מתחיל בעבודה ישירה ושיטתית, ועובר לצמיחה מבוססת קהילה ושמועה.

\section{אסטרטגיית כניסה לשוק בשלושה שלבים}\label{sec:gtm}

\begin{enumerate}
  \item \textbf{שלב 1, הוכחת היתכנות (2027)}: \RRroofsFirst{} גגות ראשונים בשכונות סמוכות בגוש דן. המטרה היא תוצאות אמיתיות, צילום תהליך ושלושה עד חמישה מסעדנים כעוגן. אין פרסום ממומן בשלב זה.
  \item \textbf{שלב 2, שכפול (2028--2029)}: התרחבות לאשכולות של גגות באותה שכונה, כדי לקצר את מסלולי האספקה. מכירה ישירה לוועדי בתים בעזרת בניינים שכבר פועלים כדוגמה, גיוס ה-\en{Series A} ומ-2029 מכירת רישיון \en{RootOS} לחוות חיצוניות.
  \item \textbf{שלב 3, קנה מידה (2030--2031)}: הגעה ל-\RRroofsFinal{} גגות, הרחבת בסיס הרישיונות של \en{RootOS} ומותג מוכר ברמת העיר.
\end{enumerate}

אנו פועלים לפי עיקרון \en{Build-Measure-Learn} של \cite{ries2011}. כל גג בשלב הראשון הוא ניסוי: מודדים מחיר שהמסעדות מוכנות לשלם, שיעור הנרשמים לסל בבניין ושיעור הנטישה, ומתקנים את ההצעה לפני הגל הבא.

\section{משפך מכירות לוועדי בתים}\label{sec:funnel}

ועד בית הוא לקוח קשה: ההחלטה קולקטיבית, נדרש רוב של בעלי הדירות, ומחזור המכירה ארוך. טבלה~\ref{tab:funnel} מציגה את משפך המכירות עם הנחות ההמרה שלנו. אלו הנחות עבודה של הצוות, שיתעדכנו אחרי השלב הראשון.

\begin{table}[htbp]
\centering\small
\rrstripes
\begin{tabularx}{\linewidth}{>{\bfseries}p{3.6cm}XXX}
\rrtablehead
\rrth{שלב במשפך} & \rrth{תיאור} & \rrth{המרה לשלב הבא} & \rrth{משך משוער} \\
בניינים מתאימים (שש קומות ומעלה) & מיפוי לפי גובה, שנת בנייה וגג פנוי & 30\% & שבועיים \\
פנייה ופגישה ראשונה עם הוועד & הצגה קצרה ודוגמה מבניין פועל & 40\% & חודש \\
סקר מבני ראשוני & מהנדס בודק עומסים וגישה & 60\% & חודש \\
אסיפת דיירים ואישור & הצגת ההסכם והחלוקה & 50\% & חודשיים \\
חוזה חתום ל-\RRcontractYears{} שנים & חתימה והזמנת התקנה & 100\% & שבועיים \\
\end{tabularx}
\caption{משפך מכירות לוועדי בתים: הנחות המרה (הערכת צוות).}
\label{tab:funnel}
\end{table}

המשמעות: מתוך כ-28 בניינים מתאימים נחתם כ-1 חוזה, כלומר המרה כוללת של כ-3.6\%. לכן חשוב למלא את ראש המשפך בבניינים בצפיפות גבוהה ולשפר את ההמרה בעזרת בניינים פועלים כדוגמה.

\section{ערוץ המסעדות}\label{sec:restaurants}

המסעדות הן ההכנסה הראשונה ביותר, והמכירה אליהן קצרה: טעימה בשבוע הראשון, חוזה אספקה לשלושה חודשים. אנו מציעים שפים ובעלי בתי קפה בתוך קילומטר מהגג. התוצרת מגיעה באותו יום, ובתפריט מופיע שם הבניין. עבור המסעדן זה יתרון שיווקי, ועבורנו ביקוש יציב ללא תלות בסוחרים.

\section{קהילת התושבים ואפליקציה}\label{sec:community}

מנוי \enquote{סל ירוק} נמכר בתוך הבניין ובבניינים הסמוכים. אפליקציית התושבים מאפשרת הזמנה, בחירת סל ומעקב אחר הגידול של הצמחים. כל בניין הוא קהילה קטנה, ולכן ההמלצה האישית היא ערוץ ההפצה העיקרי. אירוע קטיף פתוח בכל רבעון מחזק את הקשר עם הדיירים.

\section{קול המותג}\label{sec:brand}

המותג מדבר בשפה פשוטה, מקומית ושקופה: \enquote{גדל מעליך}. אנו מראים נתונים אמיתיים, כגון מתי נקטף הירק וכמה מים נחסכו, ולא סיסמאות. הטון חם ולא ראוותני.

\section{מדדי הצלחה}\label{sec:kpis}

\begin{itemize}
  \item \textbf{עלות רכישת לקוח (\en{CAC})}: עלות שיווק ומכירה ללקוח חדש, בנפרד למסעדה, לתושב ולוועד.
  \item \textbf{נטישה (\en{churn}) של מנויי הסל}: אחוז המנויים שעוזבים בחודש. יעד: נמוך מ-5\% בחודש.
  \item \textbf{שיעור המרה במשפך הוועדים}: מול ההנחות בטבלה~\ref{tab:funnel}.
  \item \textbf{ממוצע הזמנה למסעדה ושיעור חידוש חוזים}.
\end{itemize}

\begin{insightbox}[תובנה]
המכירה הקשה היא ועד הבית, והמכירה הקלה היא התוצרת. לכן אנו משקיעים בשיווק ראשוני בבניינים הראשונים, ואחר כך נשענים על ההמלצות שלהם.
\end{insightbox}
